\section{Detailed benchmark results and scientific attribution}
\label{sec:si-detailed-results}

\subsection{Complete workflow and tier comparisons}

The publication dataset contains 200 task runs: one retained run for each of 50 tasks under Guided, Standard, Record, and Brief. Each run has a composite score and all three component scores. Table~\ref{tab:si-arm-tier} reports the means without pooling away the task tiers. Standard is the actual comparison branch using the same model backbone. Record and Brief vary the workflow's persistent acceptance requirements and diagnostic feedback, respectively.

\begin{table}[!htbp]
\centering
\small
\caption{Composite and component means for each workflow and task tier.}
\label{tab:si-arm-tier}
\begin{tabular}{llrrrrr}
\toprule
Tier & Workflow & $n$ & Composite & Graph & SC & Accuracy \\
\midrule
All & Guided & 50 & 0.7803 & 0.8788 & 0.9233 & 0.6924 \\
 & Standard & 50 & 0.5940 & 0.7837 & 0.5303 & 0.5943 \\
 & Record & 50 & 0.7044 & 0.7876 & 0.7177 & 0.6839 \\
 & Brief & 50 & 0.7415 & 0.8430 & 0.8468 & 0.6720 \\
\addlinespace
Easy & Guided & 20 & 0.8809 & 0.9620 & 0.9250 & 0.8454 \\
 & Standard & 20 & 0.7305 & 0.9007 & 0.5858 & 0.7744 \\
 & Record & 20 & 0.9443 & 0.9576 & 0.9050 & 0.9618 \\
 & Brief & 20 & 0.9054 & 0.9523 & 0.9125 & 0.8940 \\
\addlinespace
Normal & Guided & 10 & 0.7370 & 0.8589 & 0.9021 & 0.6342 \\
 & Standard & 10 & 0.5731 & 0.7607 & 0.5217 & 0.5675 \\
 & Record & 10 & 0.5287 & 0.6786 & 0.5783 & 0.4789 \\
 & Brief & 10 & 0.5857 & 0.7472 & 0.7740 & 0.4646 \\
\addlinespace
Hard & Guided & 10 & 0.9577 & 0.9107 & 0.9542 & 0.9672 \\
 & Standard & 10 & 0.6887 & 0.8481 & 0.5033 & 0.7549 \\
 & Record & 10 & 0.8723 & 0.9243 & 0.6650 & 0.9672 \\
 & Brief & 10 & 0.8282 & 0.9286 & 0.7767 & 0.8372 \\
\addlinespace
Impossible & Guided & 10 & 0.4451 & 0.7006 & 0.9100 & 0.1700 \\
 & Standard & 10 & 0.2473 & 0.5083 & 0.4550 & 0.1000 \\
 & Record & 10 & 0.2325 & 0.4200 & 0.5350 & 0.0500 \\
 & Brief & 10 & 0.4829 & 0.6345 & 0.8583 & 0.2700 \\
\bottomrule
\end{tabular}
\par\smallskip\begin{minipage}{\linewidth}\footnotesize SC denotes the graded self-consistency component. Every mean is calculated from the complete publication table. The overall tier weights are 20/50 for Easy and 10/50 for each other tier.\end{minipage}
\end{table}


Table~\ref{tab:si-pairwise} gives all six pairwise comparisons on the same 50 tasks. Guided--Standard has a mean difference of 0.186302, a median difference of 0.107950, and 35 task wins, eight losses, and seven ties. The 95\% task-bootstrap interval is [0.100268, 0.275854]. Within the Normal tier, the means are 0.737020 and 0.573050, with a difference of 0.163970 [0.036520, 0.292230] and seven wins in ten tasks. Within Hard, the means are 0.957650 and 0.688750; the difference is 0.268900 [0.121300, 0.445930], with Guided higher in all ten task pairs. Easy has means of 0.880940 and 0.730485, with 11 wins, five losses, and four ties; its paired median difference is 0.020850. Impossible has means of 0.445060 and 0.247330 and is examined at the level of the requested target below.

\begin{table}[!htbp]
\centering
\small
\caption{All pairwise composite-score comparisons across the 50 tasks.}
\label{tab:si-pairwise}
\begin{tabular}{lrrr}
\toprule
Comparison & Mean difference & 95\% interval & W/L/T \\
\midrule
Guided--Record & +0.0759 & [+0.0121, +0.1438] & 27/8/15 \\
Guided--Standard & +0.1863 & [+0.1003, +0.2759] & 35/8/7 \\
Guided--Brief & +0.0388 & [-0.0317, +0.1076] & 26/10/14 \\
Record--Standard & +0.1104 & [+0.0167, +0.2062] & 28/13/9 \\
Record--Brief & -0.0371 & [-0.0999, +0.0258] & 15/23/12 \\
Standard--Brief & -0.1475 & [-0.2474, -0.0494] & 12/33/5 \\
\bottomrule
\end{tabular}
\par\smallskip\begin{minipage}{\linewidth}\footnotesize The difference is the first workflow minus the second. W/L/T gives the first workflow's task wins, losses, and ties. Percentile task-bootstrap intervals use 100,000 resamples and seed 20261031; the four configurations of each sampled task remain together.\end{minipage}
\end{table}


The component decomposition in Table~\ref{tab:si-weighted-gain} expresses each contribution in composite-score units. Across all tasks, the self-consistency contribution is +0.117875 and the accuracy contribution +0.058916. In Normal and Hard together, the composite difference is +0.216435, including +0.124688 from self-consistency and +0.083712 from accuracy. Hard's accuracy contribution, +0.127391, accompanies a +0.135250 self-consistency contribution. These quantities describe the requested comparative study at two levels: the answer reached and the evidence supporting its attribution.

\begin{table}[!htbp]
\centering
\small
\caption{Weighted decomposition of the Guided--Standard composite difference.}
\label{tab:si-weighted-gain}
\begin{tabular}{lrrrrr}
\toprule
Scope & $n$ & Composite & $0.1\Delta G$ & $0.3\Delta SC$ & $0.6\Delta A$ \\
\midrule
All & 50 & +0.186302 & +0.009513 & +0.117875 & +0.058916 \\
Easy & 20 & +0.150455 & +0.006130 & +0.101750 & +0.042579 \\
Normal & 10 & +0.163970 & +0.009823 & +0.114125 & +0.040033 \\
Hard & 10 & +0.268900 & +0.006259 & +0.135250 & +0.127391 \\
Impossible & 10 & +0.197730 & +0.019226 & +0.136500 & +0.042000 \\
Normal + Hard & 20 & +0.216435 & +0.008041 & +0.124688 & +0.083712 \\
\bottomrule
\end{tabular}
\par\smallskip\begin{minipage}{\linewidth}\footnotesize Component contributions are calculated from the recorded unrounded components. Composite differences retain released composite scores, so the final displayed digits can differ by score rounding.\end{minipage}
\end{table}


The ablations give additional resolution. Guided--Record averages +0.075892 [0.012134, 0.143818] across all tasks. In Normal and Hard together, its +0.146845 composite difference decomposes into +0.008338 from graph, +0.091938 from self-consistency, and +0.046574 from accuracy. Guided--Brief averages +0.038816 [$-0.031704$, 0.107648], with mean differences of +0.151350 on Normal and +0.129470 on Hard. The task-level records retain both the general patterns and the crossed outcomes: NB1 has composite scores of 0.9628, 0.9128, and 0.9900 for Guided, Standard, and Brief; Guided and Brief both have self-consistency 1. NI1 has Guided and Standard composites of 0.9958 and 0.6196, with self-consistency 1 and 0.35. NQ5 favors Record over Guided, while NU3 favors Brief. Multistep performance is therefore resolved through the individual pathways and their evidence relationships.

\subsection{Runtime delivery and the organized scientific record}

During execution, Guided, Record, and Brief each produced 50 root mechanism graphs. Standard produced 43. Its remaining seven studies---F1, HA1, HF1, IA9, M17, NR1, and P3---were organized offline from their own conversations, reports, tool results, and calculation artifacts. The resulting analysis contains 50 organized graphs per workflow and 200 complete score profiles. Offline organization follows the evaluation procedure described in the methods and retains the source of each recovered field.

The seven Standard studies have a mean organized composite of 0.153743, compared with 0.828529 for Guided on the same tasks. Their contribution is 50.7\% of the total Guided--Standard score difference. The scientific content retained in a run and the production of its final root artifact are both visible in these records. The self-consistency distribution in Table~\ref{tab:si-sc-distribution} further resolves how much evidence was assembled for the 40 non-Impossible tasks. Guided has 24 full-score and 16 partial-score records; Standard has seven full-score, 29 partial-score, and four zero-score records.

\begin{table}[!htbp]
\centering
\small
\caption{Self-consistency distribution over all 40 non-Impossible tasks.}
\label{tab:si-sc-distribution}
\begin{tabular}{lrrrr}
\toprule
Workflow & SC $=1$ & $0<$ SC $<1$ & SC $=0$ & Mean SC \\
\midrule
Guided & 24 & 16 & 0 & 0.9266 \\
Standard & 7 & 29 & 4 & 0.5492 \\
Record & 16 & 24 & 0 & 0.7633 \\
Brief & 20 & 20 & 0 & 0.8439 \\
\bottomrule
\end{tabular}
\par\smallskip\begin{minipage}{\linewidth}\footnotesize This table describes the graded component for the positive-chemistry tiers. The separate scientific-acceptance record is evaluated from the complete task contract.\end{minipage}
\end{table}


\subsection{Elapsed research time and accumulated calculation effort}

Table~\ref{tab:si-batch-times} preserves every batch observation from the 48-run repeated-task experiment. A batch contains the same four tasks at a concurrency of two, and its elapsed time spans the batch's actual start and finish. Guided's mean is 4.4514 h, Standard's 3.9577 h, Record's 3.8942 h, and Brief's 4.8273 h. Relative to Standard, Guided's observed mean is 12.5\% longer; relative to Brief, it is 7.8\% shorter. These observations include waiting, calculation execution, and continuation of the research after results return.

\begin{table}[!htbp]
\centering
\small
\caption{All 12 elapsed-time observations in the repeated-task study, in hours.}
\label{tab:si-batch-times}
\begin{tabular}{lrrrr}
\toprule
Workflow & Batch 1 & Batch 2 & Batch 3 & Mean \\
\midrule
Guided & 4.2383 & 3.8669 & 5.2489 & 4.4514 \\
Standard & 5.3381 & 2.5044 & 4.0306 & 3.9577 \\
Record & 5.4189 & 2.6203 & 3.6433 & 3.8942 \\
Brief & 7.2689 & 3.6475 & 3.5656 & 4.8273 \\
\bottomrule
\end{tabular}
\par\smallskip\begin{minipage}{\linewidth}\footnotesize Each observation is a four-task batch (E1, NB1, HA1, IA1) at concurrency two. Three batches per workflow contain 48 task runs in total. Batch completion times include the original research lifecycle and outcomes.\end{minipage}
\end{table}


The P2 effort comparison uses the common 15-task set D3, M1, M8, NB1, NQ5, NR1, NU3, NU5, NU6, HA1, HA6, HB1, HB10, HD2, and HF1. This yields 60 run-level effort records, with all four configurations represented for each task. Table~\ref{tab:si-common-effort} gives the total and mean job-hours, job counts, terminal outcomes, and corresponding final scores. Recorded ledger durations are summed across quantum-chemical jobs. For Standard/M1, the archived effort reconstruction uses native-log CPU time divided by four, as specified in the source record. The measurement unit throughout this comparison is accumulated job-hours.

\begin{table}[!htbp]
\centering
\small
\caption{Recorded quantum-chemical effort on the 15 tasks shared by all configurations.}
\label{tab:si-common-effort}
\begin{tabular}{lrrrrrr}
\toprule
Workflow & Job-h total & Job-h/task & Jobs & Failed & Cancelled & Mean $S$ \\
\midrule
Guided & 166.19 & 11.079 & 789 & 114 & 94 & 0.8124 \\
Standard & 113.32 & 7.555 & 419 & 45 & 50 & 0.6495 \\
Record & 118.16 & 7.877 & 431 & 62 & 38 & 0.7360 \\
Brief & 233.85 & 15.590 & 883 & 167 & 120 & 0.7273 \\
\bottomrule
\end{tabular}
\par\smallskip\begin{minipage}{\linewidth}\footnotesize Each workflow contributes 15 task runs. Job-hours sum recorded process durations. Failed and cancelled jobs are separate terminal categories; job counts do not measure scientific acceptance.\end{minipage}
\end{table}


Guided's mean of 11.0793 job-hours accompanies a mean composite of 0.812413. Relative to Standard, it uses 46.7\% more accumulated job time and gains 0.162933 in mean composite score. Relative to Brief, it uses 28.9\% less accumulated job time and gains 0.085080. On seven of the 15 paired tasks, Guided has both a higher score and lower accumulated job time than Brief. The job counts show how this expenditure was made: Guided records 789 jobs, including 114 failures and 94 cancellations; Brief records 883 jobs, including 167 failures and 120 cancellations. In the wider 63-run effort archive, NU3, NU5, NU6, and NR1 account for 163.5 of 673.82 job-hours across their 16 runs. Their trajectories contain repeated attempts to establish the required network and support examination of the concrete structures and pathways reached by those attempts.

\subsection{Scientific-attribution cases}

Table~\ref{tab:si-case-profiles} gives every score component for the cases below. Each comparison follows the chemical object to which the numerical or path evidence was attached.

\begin{table}[!htbp]
\centering
\small
\caption{Complete score profiles for the four scientific-attribution cases.}
\label{tab:si-case-profiles}
\begin{tabular}{llrrrr}
\toprule
Task & Workflow & Composite & Graph & SC & Accuracy \\
\midrule
HA8 & Guided & 0.9625 & 1.0000 & 0.8750 & 1.0000 \\
 & Standard & 0.6250 & 1.0000 & 0.3500 & 0.7000 \\
HH1 & Guided & 0.9838 & 1.0000 & 1.0000 & 0.9731 \\
 & Standard & 0.8338 & 1.0000 & 0.5000 & 0.9731 \\
HI1 & Guided & 0.9999 & 0.9992 & 1.0000 & 1.0000 \\
 & Standard & 0.7032 & 0.3325 & 0.2333 & 1.0000 \\
NU6 & Guided & 0.6719 & 1.0000 & 0.9375 & 0.4844 \\
 & Standard & 0.6892 & 0.8167 & 0.6250 & 0.7000 \\
\bottomrule
\end{tabular}
\par\smallskip\begin{minipage}{\linewidth}\footnotesize These profiles retain the distinct contributions of pathway coverage, evidence support, and agreement with the requested answer.\end{minipage}
\end{table}


\paragraph{HA8: stable reactant reference and reactive conformer.}
Guided reports 29.46 kcal/mol and Standard 20.43 kcal/mol, against a reference of 29.48 kcal/mol. The 9.03 kcal/mol difference is traced to their energy references: the lowest all-s-trans reactant minimum in Guided and a reactive cisoid conformer in Standard. The stable-species energy reference specifies the activation-energy denominator, while the IRC identifies the local structures connected by a particular saddle. Maintaining both associations allows a path reached from a reactive conformer to be interpreted relative to the specified reactant species.

\paragraph{HH1: a reference-consistent number followed by evidence closure.}
Both workflows report a binaphthyl racemization barrier of 23.48 kcal/mol; the reference is 23.01 kcal/mol. Their accuracy components are 0.973067 and 0.973064 and both graph components are 1. Guided has self-consistency 1, whereas Standard has 0.5. Their composites are consequently 0.9838 and 0.8338. Standard's archive contains candidate channels and IRC inputs; Guided's delivered evidence completes the saddle-to-IRC-to-endpoint associations required by the step. The pair locates a 0.1500 composite difference after closely matching numerical answers have already been obtained.

\paragraph{HI1: numerical agreement and coverage of the supported study.}
Guided and Standard both have accuracy approximately 1. Their graph components are 0.999176 and 0.332509, and self-consistency components are 1 and 0.233333. The corresponding composites are 0.9999 and 0.7032. Here, evidence closure and delivered coverage jointly determine how much of the requested study accompanies the numerical answer.

\paragraph{NU6: the optimized saddle and the intended second step.}
The Guided record gives $E_{\mathrm{TS}}=-247.030042392$ Ha for its concerted saddle region and $E_{\mathrm{betaine}}=-247.060060699$ Ha. Using 627.5094740631 kcal mol$^{-1}$ Ha$^{-1}$ gives
\begin{equation}
\left[E_{\mathrm{TS}}-E_{\mathrm{betaine}}\right]\,627.5094740631
=18.836772\ \mathrm{kcal\,mol^{-1}}.
\end{equation}
Standard's archived second-step ring-closure value is 18.837 kcal/mol. Guided retains the target betaine-to-product ring-closure saddle as unresolved and describes the computed energy cost in relation to the concerted saddle region. Its attempts from interpolated betaine geometries and constrained-scan peaks repeatedly return toward the direct reactant-to-product saddle region. Following the optimized structure and its actual endpoints is therefore necessary to assign the energy difference to the step that was established.

\subsection{Impossible-tier task and claim resolution}

Table~\ref{tab:si-imp-pairs} resolves Brief's mean composite of 0.482940 and Guided's 0.445060 into all ten task differences. Brief is higher on three tasks, Guided on five, and two tie. The Brief-minus-Guided mean is +0.037880, with the existing task-bootstrap interval [$-0.188220$, 0.283970]. Its weighted contributions are $-0.006607$ from graph, $-0.015500$ from self-consistency, and +0.060000 from accuracy. IA1 and IA9 provide the two large positive differences, while IG8 has the opposite accuracy ordering. Across the seven tasks with equal accuracy components in these two workflows, Brief-minus-Guided averages $-0.050000$.

\begin{table}[!htbp]
\centering
\small
\caption{Task-resolved comparison of Brief and Guided in the Impossible tier.}
\label{tab:si-imp-pairs}
\begin{tabular}{lrrrrr}
\toprule
Task & Guided C & Brief C & Brief $-$ Guided & Guided A & Brief A \\
\midrule
IA1 & 0.1717 & 1.0000 & +0.8283 & 0.0 & 1.0 \\
IA5 & 0.4000 & 0.3667 & -0.0333 & 0.0 & 0.0 \\
IA7 & 0.2583 & 0.2583 & +0.0000 & 0.0 & 0.0 \\
IA9 & 0.3667 & 0.9333 & +0.5666 & 0.0 & 1.0 \\
IB9 & 0.8200 & 0.7700 & -0.0500 & 0.7 & 0.7 \\
IC3 & 0.4000 & 0.4000 & +0.0000 & 0.0 & 0.0 \\
ID2 & 0.3339 & 0.1839 & -0.1500 & 0.0 & 0.0 \\
IE2 & 0.3333 & 0.1833 & -0.1500 & 0.0 & 0.0 \\
IF3 & 0.3667 & 0.4000 & +0.0333 & 0.0 & 0.0 \\
IG8 & 1.0000 & 0.3339 & -0.6661 & 1.0 & 0.0 \\
\bottomrule
\end{tabular}
\par\smallskip\begin{minipage}{\linewidth}\footnotesize C denotes the composite and A the accuracy component for the task-defined numerical or negative target. All values are retained from the publication score table.\end{minipage}
\end{table}


Brief reaches accuracy 1 for IA1, the square-geometry cyclobutadiene automerization target, and IA9, dioxetane thermolysis. Guided reaches accuracy 1 for IG8, the chromium-dimer target. Both have accuracy 0.7 in IB9, which contains positive and negative target components. Across all 40 Impossible runs, four records have accuracy 1 and 14 have self-consistency 1 with accuracy 0. These component combinations are reported together with the identity of the requested state, channel, or property.

The preserved Brief/IC3 graph identifies the requested lowest open-shell singlet $^1A_2$ as inaccessible under the permitted route, marks its state-change edge accordingly, and withholds the requested singlet--triplet gap. It retains a separately characterized closed-shell $^1A_1$ branch. The PM7 value of +45.9 kcal/mol and GFN2-xTB value of $-1.7$ kcal/mol are described as method-dependent exploratory relative energies and are explicitly excluded from the requested gap. The graph records that the restricted calculation with a mixed initial guess does not test the requested unrestricted electronic state. The released profile remains Graph/SC/Accuracy = 1/1/0 and composite 0.4.

The Record/IA5 archive similarly distinguishes the requested disrotatory ring opening from an allowed conrotatory path. It marks the requested disrotatory edge inaccessible and records the positive conrotatory path separately at 33.57 kcal/mol. The archived adjudication documents that the matcher selected the positive path sharing the endpoints, leaving the accuracy component at 0. These records preserve both the delivered scientific claims and the values used in the publication analysis.

\subsection{Easy-tier energy references and delivery trajectories}
\label{sec:si-easy-case-audit}

The five Easy tasks examined here distinguish numerical reference choice, chemical identity matching, and the delivery of computed evidence. We inspected the retained runs' final mechanism graphs, native calculation logs, evaluation notes, and runtime events. Table~\ref{tab:si-easy-case-audit} preserves the released scores. $\mathrm{SC}_{\mathrm{all}}$ grades the delivered graph, whereas $\mathrm{SC}_{\mathrm{Gold}}$ grades its reference-matched subgraph and excludes additional elements. Their difference identifies a change in the objects being assessed; the composite uses $\mathrm{SC}_{\mathrm{all}}$.

\begin{table}[!htbp]
\centering
\small
\setlength{\tabcolsep}{4pt}
\caption{Released component scores for the five inspected Easy tasks, rounded to four decimal places. Brief denotes Brief-errors. The two SC columns retain their separate assessment scopes. Accuracy is the final component after reference matching and the specified numerical penalty cap.}
\label{tab:si-easy-case-audit}
\begin{tabular}{llrrrrr}
\toprule
Task & Workflow & Graph & $\mathrm{SC}_{\mathrm{all}}$ & $\mathrm{SC}_{\mathrm{Gold}}$ & Accuracy & Composite \\
\midrule
F4 & Guided & 0.6009 & 1.0000 & 0.0000 & 0.0000 & 0.3601 \\
 & Standard & 0.8333 & 0.7500 & 0.7500 & 1.0000 & 0.9083 \\
 & Record & 0.9167 & 0.7500 & 0.7500 & 1.0000 & 0.9167 \\
 & Brief & 0.6009 & 1.0000 & 0.0000 & 0.0000 & 0.3601 \\
\addlinespace
G2 & Guided & 0.8889 & 0.5000 & 0.5000 & 0.0000 & 0.2389 \\
 & Standard & 1.0000 & 0.5000 & 0.5000 & 0.8782 & 0.7769 \\
 & Record & 1.0000 & 0.5000 & 0.5000 & 0.8782 & 0.7769 \\
 & Brief & 1.0000 & 0.5000 & 0.5000 & 1.0000 & 0.8500 \\
\addlinespace
D1 & Guided & 1.0000 & 1.0000 & 1.0000 & 0.7000 & 0.8200 \\
 & Standard & 0.7917 & 0.5000 & 0.5000 & 1.0000 & 0.8291 \\
 & Record & 0.7917 & 1.0000 & 1.0000 & 0.7000 & 0.7992 \\
 & Brief & 1.0000 & 0.7500 & 0.7500 & 0.7000 & 0.7450 \\
\addlinespace
D3 & Guided & 1.0000 & 0.7500 & 0.7500 & 0.7000 & 0.7450 \\
 & Standard & 1.0000 & 0.7500 & 0.7500 & 0.9995 & 0.9247 \\
 & Record & 0.9444 & 0.3500 & 0.3500 & 0.9909 & 0.7940 \\
 & Brief & 0.8333 & 0.7500 & 0.7500 & 0.7000 & 0.7283 \\
\addlinespace
L2 & Guided & 1.0000 & 0.7500 & 0.7500 & 0.7000 & 0.7450 \\
 & Standard & 1.0000 & 0.5000 & 0.5000 & 0.7000 & 0.6700 \\
 & Record & 1.0000 & 0.7500 & 0.7500 & 0.8832 & 0.8549 \\
 & Brief & 1.0000 & 0.8333 & 0.7500 & 0.8832 & 0.8799 \\
\bottomrule
\end{tabular}
\end{table}

\paragraph{F4: reference identity and a verified adduct pathway.}
Guided delivered hydride transfer from the Lewis adduct of borane and formaldehyde to methoxyborane. Its bound reactant and saddle energies were $-141.041124115$ and $-141.015601446$~hartree, giving 16.01572~kcal~mol$^{-1}$ and supporting the declared 16.02. Independent saddle frequency analysis gave one imaginary frequency, $-398.7218$~cm$^{-1}$. Both actual IRC outer frames subsequently converged to the corresponding zero-imaginary-frequency endpoint minima. The four workflows obtained essentially the same saddle and adduct energies.

The campaign reference at revision \texttt{53309d0e} encoded the reactant as \texttt{[BH3].C=O}, although its step description and 16.02-kcal~mol$^{-1}$ reference concerned the Lewis adduct. Guided and Brief instead declared the bonded adduct. The original Guided evaluation explicitly reported that no delivered step matched the reference endpoint species. Revision \texttt{0b7b5e5b} corrected the reference node to \texttt{[BH3-][O+]=C}; the retained offline protocol used the earlier campaign taskbank. The released Standard and Record scores include their recorded annotation-identity recovery. F4 therefore contains a documented reference-identity effect: Guided's $\mathrm{SC}_{\mathrm{all}}=1$ coexists with $\mathrm{SC}_{\mathrm{Gold}}=0$ and Accuracy~0 despite the archived 16.02 pathway. Table~\ref{tab:si-easy-case-audit} retains this publication record.

\paragraph{G2: computed evidence and final structured delivery.}
Guided's native logs contain a concerted double-H-transfer saddle at $-189.093172517$~hartree with one imaginary frequency, $-444.9504$~cm$^{-1}$. The optimized cis reactant complex and the IRC-derived reactant minimum have energies of $-189.094518484$ and $-189.094518735$~hartree. They give barriers of 0.84461 and 0.84476~kcal~mol$^{-1}$, respectively, close to the 0.87 reference.

The runtime recorded four rejections of a numerical barrier on a companion complex-product edge whose scientific checks remained open. Meanwhile, delivery reports for the selected separated-products target returned \texttt{complete=true} with no unresolved target. After an initial deletion, a Supervisor request prompted restoration of all numerical barriers and the saddle's energy/log fields. The fourth barrier rejection was followed by deletion of the companion-edge values. A subsequent artifact check rejected the saddle-evidence claim; the next recorded edit removed the selected edge's remaining numerical barrier and the saddle's energy/log fields. The successful write and program output confirm these removals. The final graph retained the values in descriptive text, so the numerical target received Accuracy~0. Its graded SC contribution remained 0.5.

Brief retained a numerical barrier of 0.844 and obtained Accuracy~0.99999. Standard and Record declared approximately 3.4 using a lower-energy reactant reference; Record's native report identifies this as a trans-diimide stack and separately records the approximately 0.84 cis-reference result.

\paragraph{D1 and D3: numerical penalties with different physical origins.}
The evaluation notes explicitly attribute Accuracy~0.7 to the maximum numerical-deviation penalty of 0.3. For D1, 15.738 versus 18.830~kcal~mol$^{-1}$ gives an underlying Gaussian score of 0.5880 before this cap. Guided and Standard obtained nearly identical saddle energies, but their reactant energies differ by 3.11859~kcal~mol$^{-1}$. Their butadiene dihedrals are approximately $26.33^\circ$ and $-179.93^\circ$, corresponding to reactive s-cis and lower s-trans references. Record also delivered the s-cis barrier and retained SC~1 with Accuracy~0.7.

D3 has a different origin. Guided and Standard delivered barriers of 18.14687 and 14.55472~kcal~mol$^{-1}$. Their reactant energies differ by only 0.09899~kcal~mol$^{-1}$, whereas their saddle energies differ by 3.69113~kcal~mol$^{-1}$. Both saddle logs contain exactly one imaginary frequency. The archived evaluations separately assign SC~0.75 with product-endpoint association warnings; the Accuracy cap follows the numerical barrier comparison.

\paragraph{L2: two reference barriers and the selected comparison edge.}
Guided used the reactive cZc minimum and obtained 18.98740~kcal~mol$^{-1}$; its graph states that the stable-to-reactive conformer contribution was uncomputed. Record and Brief used the lower tZt reference, giving approximately 28.275. The conformer energy difference is 9.28761~kcal~mol$^{-1}$.

Standard delivered both references and both barriers: 18.99 on edge \texttt{e1} and 28.28 on \texttt{e2}. The archived evaluation selected \texttt{e1}. Both candidates have identical bond-change lists, and the unspecified-stereochemistry reference gives equal endpoint substitution costs for these reactants. The final edge-identifier tie-break selects \texttt{e1}, so the local-reference value receives the capped Accuracy~0.7. Standard's final report explicitly contains the stable-reference barrier and conformer gap.

\subsection{Sensitivity to target-object assignment in Impossible tasks}
\label{sec:matcher-sensitivity}

We examined two cases in which distinct scientific objects share the same molecular-connectivity endpoints. The analysis uses the original final mechanism graphs and process records for Record/IA5 and Brief/IC3. It asks how the Impossible-tier means change if the existing negative-credit rule is applied to the object requested by the task. The 200 published scores remain unchanged. For this conditional calculation, we hold Graph and Self-consistency fixed in both cases and retain all scores and components for the other 198 runs.

IA5 requests the thermal disrotatory ring opening of cyclobutene. Record's final graph marks the disrotatory edge \texttt{e\_dis} and its transition-state node inaccessible and declares no barrier for that pathway. A separate conrotatory edge, \texttt{e\_con}, carries a barrier of 33.57 kcal/mol. The final process record explicitly distinguishes the requested disrotatory target from this computed alternative. Both edges have the same connectivity endpoints and bond-change signature, so these features alone do not identify which pathway answers the task.

IC3 requests the lowest singlet state of phenylnitrene and its singlet--triplet gap. Brief's original graph distinguishes the closed-shell singlet \texttt{S1\_1A1} from the open-shell target \texttt{S2\_1A2}. It marks the latter node and its state-change edge inaccessible, withholds the requested numerical gap, and explicitly excludes the closed-shell branch's relative energies from that gap. The two state-change edges again have identical endpoint SMILES and bond-change signatures. The published baseline for Brief/IC3 is 0.4 after the documented evidence-reference formatting repair; its raw diagnostic score of $-1$ is not used here.

In both original graphs, the current matcher selects the positive alternative object. The intended negative target passes the unchanged negative-claim predicate when assigned to the corresponding reference target. A correction must therefore resolve pathway or electronic-state identity before evaluating the final claim. Simply preferring a negative edge would allow a refusal to conceal a contradictory positive claim about the same object. The positive alternative should be excluded only when the run's own final statements establish that it is a different scientific object.

Both published cases have Graph $=1$, Self-consistency $=1$, Accuracy $=0$, and composite $=0.4$. Assigning Accuracy $=1$ to either intended negative target changes that case's conditional composite by $0.6$, following the fixed Accuracy weight. Its configuration's ten-task Impossible mean therefore increases by $0.06000$. Table~\ref{tab:matcher-sensitivity} reports the four scenarios; the accompanying \texttt{matcher\_sensitivity\_scenarios.csv} retains their configuration-level means and ranks.

\begin{table}[htbp]
\centering
\small
\caption{Impossible-tier mean composite scores under two conditional target-assignment corrections. Only the indicated Accuracy values change; Graph, Self-consistency, and the other 198 runs remain fixed. The original row is the published baseline.}
\label{tab:matcher-sensitivity}
\begin{tabular}{@{}lrrrr@{}}
\toprule
Scenario & Guided & Standard & Record & Brief \\
\midrule
Original published values & 0.44506 & 0.24733 & 0.23250 & 0.48294 \\
Record/IA5 target credit & 0.44506 & 0.24733 & 0.29250 & 0.48294 \\
Brief/IC3 target credit & 0.44506 & 0.24733 & 0.23250 & 0.54294 \\
Both target credits & 0.44506 & 0.24733 & 0.29250 & 0.54294 \\
\bottomrule
\end{tabular}
\end{table}

Across these scenarios, the Impossible-tier ordering Brief $>$ Guided $>$ Record is stable. Record moves above Standard when its IA5 target receives negative credit; Brief remains first when its IC3 target receives credit. Thus the two assignments can reverse the Record--Standard ordering, while leaving every Guided--Standard score, paired difference, and primary comparison unchanged. These are conditional arithmetic results for the two identified assignments, not an adopted revision of the publication scores or a full rerun of scientific validation.


\subsection{Discussion of the supplementary comparisons}

The Normal and Hard comparisons resolve the framework's contribution through relations that must survive across many local decisions: species and conformer identity, reaction-channel identity, the reference defining a reported energy, and the association of each result with the claim it supports. Generating a plausible next action from a current language context does not itself maintain this global state. A persistent evidence graph makes these accumulated relations available for subsequent reasoning and revision, complementing work on external planning and tool feedback.\citep{valmeekam2023planning,ref18,huang2024selfcorrect} The task cases further distinguish repairing an incomplete local proof from reconsidering which scientific question the current calculation can answer.

P2 contains one retained run per task and configuration. Its task-cluster bootstrap intervals describe variation across the sampled task mix, not stochastic variation from independent runs of the same task. The separate P1b study supplies three repeats on each of four tasks under each configuration; these 48 runs do not provide repeated measurements for every task in the 50-task P2 study. The complete organized P2 table covers all 200 runs and their three score components. Some Hard-task accuracy targets concern qualitative channel identity, so the reported gain includes pathway identification and evidence completion as well as numerical barriers. Repeating the full task set would characterize run-to-run variability and the frequency with which each research route is selected. The Impossible rubric targets the specified allowed-method contract; its negative-target self-consistency convention is distinct from a positive TS--IRC--endpoint chain. IA5 and IC3 show why score interpretation also requires channel and electronic-state identities. Their released scores remain unchanged. The Brief comparison is consistent with task-dependent effects of diagnostic content, while the final outcomes alone do not identify a causal effect of diagnostic length. The NU6 energy fingerprint identifies an assignment question; native geometry and pathway evidence are the independent basis for establishing whether two optimized structures represent the same saddle. Timing results retain their measured units: pilot batch elapsed time and P2 accumulated job-hours. Separating queueing, agent-active time, and end-to-end P2 latency is a further experimental question.


Cloudra's default preorganization samples 32 rigid relative placements, applies local MMFF94 relaxation, and selects candidates by energy and clustering. This is a tractable candidate-generation procedure, not an exhaustive search for the global minimum of every reactant arrangement or intermediate. CREST is available through a separate optional backend; the 200-run study does not establish uniform CREST coverage. Highly flexible systems may require broader sampling and more quantum-chemical calculations. Further work should measure conformer coverage and improve the allocation of calculations among competing candidates.

The benchmark compares electronic energy differences. It does not add zero-point, thermal/entropic, or standard-state corrections, and it does not quantify the free-energy contribution of solvation. This scope matters for NI5, specified in weakly acidic aqueous solution, and NU5, specified in dichloromethane at high concentration. Their scores assess the prescribed electronic-energy references rather than experimental solution-phase free-energy barriers. No solvent-sensitivity comparison was performed, so the present results do not determine the magnitude or direction of a solvent effect.

The atom-mapping coefficients and ambiguity window, and the decay, depth, and congestion settings in the branch-crowding descriptor, were fixed across the workflow comparisons. We did not vary them in a parameter sweep or measure each setting's separate contribution. The imaginary-mode alignment ratio of 0.5 is advisory; the stationary-point and IRC/endpoint evidence determine scientific acceptance. All P2 configurations used DeepSeek-v4.1-flash (1017), with the API alias \texttt{deepseek-v4-flash} and the recorded request settings. Comparisons with additional backbones would test how the workflow gains carry across models.


The Easy-tier trajectories distinguish reference identity from delivery policy. F4 preserves a physically checked adduct pathway but loses reference-matching credit. G2 preserves computed values in prose after validation feedback prompts removal of their structured fields. Preserving a provisional quantity and its source separately from its acceptance status would allow scientific review to restrict its interpretation without erasing the observation. The conditional IA5/IC3 analysis likewise changes target assignment while retaining the scientific-acceptance rules and the original publication dataset.
